% !PRINTED 3
\weberpage{3}

If $a$ and $a_1$ are relatively prime numbers, then the last divisor is
$a_\nu=1$. Multiplying, under this assumption, the preceding equations
by an arbitrary number $b$, it follows that $b$ is the greatest common
divisor of $ab$ and $a_1b$, and hence that every common divisor of $ab$
and $a_1$, or of $a$ and $a_1b$, must be a divisor of $b$. Thus if $a_1$
is relatively prime to $a$ and to $b$, it is also relatively prime to
$ab$; and from the special assumption that $a_1$ is a prime number it
follows that a product can be divisible by a prime number only when at
least one of the factors is divisible by it.

Thus, if the product $ab$ is divisible by $a_1$ and $a_1$ is relatively
prime to $a$, then $b$ must be divisible by $a_1$.

Let $a,b$ be any two numbers with greatest common divisor $d$, and let
$a=da'$, $b=db'$. Then $a'$ and $b'$ are relatively prime to one
another. Every number $m$ which is at once a multiple of $a$ and of $b$
has the form $m=am'=da'm'$, and in this $a'm'$ must be a multiple of
$b'$; therefore $m'$ too must be a multiple of $b'$. That is, every
number which is at once a multiple of $a$ and of $b$ is divisible by
$a'b'd$. This number $a'b'd$, which itself is a common multiple of $a$
and $b$, is therefore called the \gesperrt{least common multiple} of
$a$ and $b$.

If a number is divisible by two relatively prime numbers, then it is
also divisible by their product; and if, therefore, a number $m$ is
divisible by several numbers of which each two are relatively prime to
one another, then it is also divisible by the product of all these
numbers.

On this rests the proof of the important theorem that a number $m$ can
always be represented, and in only one way, as a product of prime
numbers. For since a divisor cannot be greater than the dividend, $m$
can certainly be divisible by only a finite number of prime numbers. If
$a$ is one of these prime numbers and $a^\alpha$ the highest power of
$a$ which divides $m$, then $\alpha$ too is a definite number. Let
likewise $b^\beta,c^\gamma,\ldots$ be the highest powers of the remaining
prime numbers occurring in $m$, namely $b,c,\ldots$, by which $m$ can
be divided. Then, since any two of the
numbers $a^\alpha,b^\beta,c^\gamma,\ldots$ are relatively prime, $m$
must also be divisible by the product
$a^\alpha b^\beta c^\gamma\cdots$; and $m$ must be equal to this product,
since otherwise another
%CONT>>
